Fix \(S\subseteq\mathcal U\) and an arm \(\rho\). Define the closed polytope
\[
F_{S,\rho}=\left\{q:\ q\geq0,\ A_Dq=m_D,\ R_eq=p_e\ (e\in S),
\ (n_{\rho,D}-n_{\sigma,D})q\geq0\ \text{for every }\sigma\in\mathcal R_{\mathrm{arm}}\right\}.\tag{1}
\]
On a point with \(d_Dq>0\), the common positive denominator makes the last inequalities equivalent to
\(Q_\rho(q)\geq Q_\sigma(q)\) for all \(\sigma\). Therefore
\[
\rho\in\operatorname{PO}(S)\quad\Longleftrightarrow\quad
F_{S,\rho}\ne\varnothing\ \text{ and }
\max_{q\in F_{S,\rho}}d_Dq>0.\tag{2}
\]
The maximum in (2) is a finite linear program; if the feasible set is empty the arm is rejected, and if its maximum is zero every candidate witness has undefined common evidence. Thus (1)--(2) compute \(\operatorname{PO}(S)\) exactly. There are finitely many subsets \(S\), so enumerating them and comparing \(\sum_{e\in S}c_e\) over those satisfying
\(\operatorname{PO}(S)=\mathcal P_{\mathrm{full}}\) characterizes every minimum-cost solution. Rational or effectively real-algebraic input makes every feasibility, optimum-sign, and cost comparison exact.

For non-submodularity, take the one-domain binary one-district model with the four response types
\[
\omega_1=(0,0),\quad\omega_2=(0,1),\quad
\omega_3=(1,0),\quad\omega_4=(1,1),\tag{3}
\]
realized by the deterministic maps \(f_{00}(t)=0\), \(f_{01}(t)=t\),
\(f_{10}(t)=1-t\), and \(f_{11}(t)=1\). Let evidence be sure, let
\(B_\rho\) be the complement of \(\{\omega_4\}\), and let \(B_\sigma=\{\omega_4\}\). Let the two genuine two-world experiment cells have rows
\[
r_1=(1,0,0,1),\qquad r_2=(0,1,1,1),\tag{4}
\]
corresponding respectively to \(\{Y_0=Y_1\}\) and
\(\{(Y_0,Y_1)\ne(0,0)\}\), and reveal value one at unit cost. These are realizable partial-copy cells, and their joint supplied values are realized by \(q=\delta_{\omega_4}\).

With normalization as the base equation, (4) gives
\[
\begin{aligned}
\mathcal K_{\{e_1\}}&=\operatorname{conv}\{\delta_{\omega_1},\delta_{\omega_4}\},\\
\mathcal K_{\{e_2\}}&=\operatorname{conv}\{\delta_{\omega_2},\delta_{\omega_3},\delta_{\omega_4}\},\\
\mathcal K_{\{e_1,e_2\}}&=\{\delta_{\omega_4}\}.
\end{aligned}\tag{5}
\]
At each of \(\delta_{\omega_1},\delta_{\omega_2},\delta_{\omega_3}\), arm \(\rho\) is uniquely optimal; at \(\delta_{\omega_4}\), arm \(\sigma\) is uniquely optimal. Hence
\[
\operatorname{PO}(\varnothing)=\operatorname{PO}(\{e_1\})
=\operatorname{PO}(\{e_2\})=\{\rho,\sigma\},\qquad
\operatorname{PO}(\{e_1,e_2\})=\{\sigma\}=\mathcal P_{\mathrm{full}}.\tag{6}
\]
Thus the marginal gain of \(e_2\) is zero at \(\varnothing\) and one after \(e_1\), violating diminishing returns. Both experiments are necessary, so the exact binary acquisition problem is
\[
\min_{x_1,x_2\in\{0,1\}}x_1+x_2\quad\text{subject to}\quad x_1+x_2\geq2.\tag{7}
\]
Its unique feasible optimum is \((1,1)\) with cost two. For the continuous relaxation, multiplying the covering inequality by \(y=1\) gives the lower bound two, equal to that feasible cost; this is the matching dual certificate.